% =====================================================================
%  Chapitre 3 -- généré par outils/convertir.py depuis docs/tp1_donnees_structurees.tex
% =====================================================================
\chapter[TP 1 -- Données structurées]{TP 1 -- Données structurées}\label{ch:tp1}
\soustitre{De la base de données au tableau de bord, en passant par Bronze, Silver et Gold}

\begin{objectifs}
  \item ingérer dans le Lakehouse les tables d'une base PostgreSQL et des fichiers CSV ;
  \item nettoyer et contrôler les données en couche Silver, avec une quarantaine ;
  \item construire une couche Gold et répondre en SQL aux questions de la direction ;
  \item charger les nouvelles données de façon incrémentale, avec \cmd{MERGE}, et relire une version passée d'une table ;
  \item entraîner un modèle qui prédit les retards de livraison, en évitant la fuite de données ;
  \item comparer Iceberg et Delta Lake, et présenter les résultats dans un tableau de bord.
\end{objectifs}

\begin{center}\small\renewcommand{\arraystretch}{1.2}
\begin{tabularx}{\textwidth}{@{}>{\bfseries}l X@{}}
\toprule
Durée indicative & 2 séances de 3 h, et environ 2 h 30 de travail personnel \\
Prérequis & chapitre~\ref{ch:environnement} : environnement installé, \og ENVIRONNEMENT PRÊT \fg{} affiché par \cmd{check\_env.py} \\
Notebooks & \nb{notebooks/tp1/}, cinq notebooks numérotés de 01 à 05 \\
\bottomrule
\end{tabularx}
\end{center}
% =====================================================================
\section{Le scénario}
% =====================================================================

Vous rejoignez l'équipe data d'\textbf{Olist}, une place de marché brésilienne qui permet à de petits
commerçants de vendre sur les grands sites de e-commerce. La direction vous pose quatre questions :

\begin{enumerate}[leftmargin=*,nosep]
  \item Comment évoluent nos ventes, mois par mois ?
  \item Quelles catégories de produits rapportent le plus ?
  \item Tenons-nous nos promesses de délai de livraison, et où avons-nous des problèmes ?
  \item Peut-on prévoir un retard de livraison dès la commande ?
\end{enumerate}

Les données sont réparties entre deux systèmes : l'application de vente, qui écrit dans une base
\textbf{PostgreSQL}, et des \textbf{fichiers CSV} exportés par d'autres services (clients, produits,
vendeurs). Votre mission est de les réunir dans le Lakehouse, de les rendre fiables, puis de répondre
aux questions.

% =====================================================================
\section{Les données}
% =====================================================================

Le jeu de données public \textbf{Olist} (Kaggle, licence CC BY-NC-SA 4.0) décrit environ
100\,000 commandes réelles passées entre 2016 et 2018. Il est fourni dans le dossier
\cmd{data/olist/} du dépôt.

\begin{center}
\small
\renewcommand{\arraystretch}{1.25}
\begin{tabularx}{\textwidth}{@{}>{\ttfamily}l l X r@{}}
\toprule
\normalfont\textbf{Table} & \textbf{Arrive par} & \textbf{Contenu} & \textbf{Lignes} \\
\midrule
orders          & PostgreSQL & Une ligne par commande : statut, dates (achat, livraison, estimation) & 99\,441 \\
order\_items    & PostgreSQL & Les articles de chaque commande : produit, vendeur, prix, frais de port & 112\,650 \\
order\_payments & PostgreSQL & Les paiements : moyen, nombre d'échéances, montant & 103\,886 \\
customers       & CSV & Les clients : identifiant, ville, état & 99\,441 \\
products        & CSV & Les produits : catégorie (en portugais), poids, dimensions & 32\,951 \\
sellers         & CSV & Les vendeurs : ville, état & 3\,095 \\
category\_translation & CSV & Traduction des catégories en anglais & 71 \\
\bottomrule
\end{tabularx}
\end{center}

\begin{attention}
dans Olist, un \cmd{customer\_id} est créé \textbf{à chaque commande}. Pour identifier une même
personne d'une commande à l'autre, il faut utiliser \cmd{customer\_unique\_id}.
\end{attention}

La figure~\ref{c3:fig:modele} montre comment les tables se relient. La table \cmd{orders} est au
centre. Une commande a plusieurs articles et plusieurs paiements, et chaque article renvoie à un
produit et à un vendeur.

\begin{figure}[H]
\centering
\resizebox{\textwidth}{!}{%
\begin{tikzpicture}[
  font=\sffamily\small, >=Stealth,
  tab/.style={draw=#1,thick,rounded corners=3pt,fill=#1!7,text width=3.6cm,align=left,
              inner sep=5pt,font=\sffamily\footnotesize},
  card/.style={font=\sffamily\scriptsize,text=black!60}]
\node[tab=gray]   (cus) at (-5.0,0)   {\textbf{customers}\\[2pt] \underline{\texttt{customer\_id}}\\ \texttt{customer\_unique\_id}\\ \texttt{customer\_state}};
\node[tab=lhblue] (ord) at (0,0)      {\textbf{orders}\\[2pt] \underline{\texttt{order\_id}}\\ \texttt{customer\_id}\\ \texttt{order\_status}\\ dates d'achat, de livraison\\ et date estimée};
\node[tab=lhblue] (itm) at (5.0,0)    {\textbf{order\_items}\\[2pt] \underline{\texttt{order\_id}, \texttt{order\_item\_id}}\\ \texttt{product\_id}, \texttt{seller\_id}\\ \texttt{price}, \texttt{freight\_value}};
\node[tab=lhblue] (pay) at (0,-3.4)   {\textbf{order\_payments}\\[2pt] \underline{\texttt{order\_id}, \texttt{payment\_seq.}}\\ \texttt{payment\_type}\\ \texttt{payment\_value}};
\node[tab=gray]   (prd) at (10.0,1.45) {\textbf{products}\\[2pt] \underline{\texttt{product\_id}}\\ \texttt{product\_category\_name}\\ poids, dimensions};
\node[tab=gray]   (sel) at (10.0,-1.8) {\textbf{sellers}\\[2pt] \underline{\texttt{seller\_id}}\\ \texttt{seller\_state}};
\node[tab=gray]   (tra) at (10.0,4.2)  {\textbf{category\_translation}\\[2pt] nom portugais $\rightarrow$ anglais};

\draw[-] (cus) -- node[card,above]{1 -- n} (ord);
\draw[-] (ord) -- node[card,above]{1 -- n} (itm);
\draw[-] (ord) -- node[card,right]{1 -- n} (pay);
\draw[-] (itm.east) -- node[card,above=2pt]{n -- 1} (prd.west);
\draw[-] (itm.east) -- node[card,below=2pt]{n -- 1} (sel.west);
\draw[-] (prd) -- (tra);

\node[font=\sffamily\scriptsize,text=lhblue,anchor=west] at (-6.9,-3.0) {\rule{8pt}{8pt}\ PostgreSQL};
\node[font=\sffamily\scriptsize,text=gray,anchor=west]  at (-6.9,-3.5) {\rule{8pt}{8pt}\ fichiers CSV};
\end{tikzpicture}}
\caption{Modèle de données Olist. Les colonnes soulignées forment la clé de chaque table, et
\og 1 -- n \fg{} se lit \og une commande, plusieurs articles \fg{}.}
\label{c3:fig:modele}
\end{figure}

% =====================================================================
\section{Ce que vous allez construire}
% =====================================================================

La figure~\ref{c3:fig:pipeline} montre le pipeline complet. Les données partent des deux sources,
traversent les couches du Lakehouse, puis servent aux usages. Les numéros renvoient aux étapes
décrites ci-dessous.

\begin{figure}[H]
\centering
\resizebox{\textwidth}{!}{%
\begin{tikzpicture}[
  font=\sffamily\small, >=Stealth,
  barre/.style={fill=lhfonce,text=white,font=\sffamily\bfseries\small,minimum height=0.6cm,inner sep=0pt},
  cadre/.style={draw=lhfonce,line width=0.8pt},
  bloc/.style={draw=lhfonce!60,line width=0.6pt,fill=white},
  et/.style={font=\sffamily\scriptsize,text=lhfonce,align=center},
  flux/.style={->,very thick,lhfonce}]

% ===== Sources =====
\draw[cadre] (0,0) rectangle (3.2,-10.4);
\node[barre,minimum width=3.2cm,anchor=north west] at (0,0) {Sources};
\node at (1.6,-2.3) {\logo[2.2em]{postgres}};  \node[et] at (1.6,-3.2) {Base de vente\\(PostgreSQL)};
\pic at (1.6,-4.6) {fichier=CSV};             \node[et] at (1.6,-5.45) {Fichiers CSV\\(référentiels)};

% ===== Plateforme Lakehouse =====
\draw[draw=lhblue,line width=2.4pt] (3.9,0) rectangle (15.9,-10.4);
\node[fill=lhblue,text=white,font=\sffamily\bfseries,minimum width=12cm,minimum height=0.7cm,anchor=north west]
  at (3.9,0) {Plateforme Lakehouse};
\draw[bloc] (4.35,-1.0) rectangle (15.2,-2.3);
\node[et,font=\sffamily\bfseries\scriptsize,anchor=west] at (4.6,-1.65) {Traitement};
\node[et,font=\sffamily\scriptsize] at (9.78,-1.65) {\logo[1.6em]{jupyter}JupyterLab\qquad\logo[1.6em]{spark}Spark (Python et SQL)};

% Tables par couche
\pic at (5.3,-4.6)  {table=gray};      \node[et] at (5.3,-5.55)  {\textbf{landing}\\fichiers bruts};
\pic at (8.1,-4.6)  {table=lhbronze};  \node[et] at (8.1,-5.55)  {\textbf{bronze}\\7 tables};
\pic at (10.9,-4.6) {table=lhsilver};  \node[et] at (10.9,-5.55) {\textbf{silver}\\tables propres};
\pic at (13.7,-4.6) {table=lhgold};    \node[et] at (13.7,-5.55) {\textbf{gold}\\prêtes à l'usage};
\pic at (10.9,-6.75) {table=gray};     \node[et,anchor=west] at (11.45,-6.75) {\textbf{quarantaine}\\lignes rejetées};

% Lakehouse ouvert : catalogue et stockage
\node[fill=lhblue,text=white,font=\sffamily\bfseries\small,minimum width=10.85cm,minimum height=0.6cm]
  at (9.78,-7.75) {Lakehouse ouvert};
\draw[bloc] (4.35,-8.2) rectangle (7.85,-10.05);
\node[barre,minimum width=3.5cm,anchor=north west] at (4.35,-8.2) {Catalogue};
\node[et] at (6.1,-9.3) {\logo[1.8em]{iceberg}\\[2pt]Iceberg REST};
\draw[bloc] (8.2,-8.2) rectangle (15.2,-10.05);
\node[barre,minimum width=7.0cm,anchor=north west] at (8.2,-8.2) {Stockage ouvert (S3)};
\node[et] at (9.3,-9.3)  {\logo[1.8em]{rustfs}\\[2pt]RustFS};
\node[et] at (11.0,-9.3) {\logo[1.8em]{iceberg}\\[2pt]Iceberg};
\node[et] at (12.6,-9.3) {\logo[1.8em]{parquet}\\[2pt]Parquet};
\node[et] at (14.2,-9.3) {\logo[1.8em]{delta}\\[2pt]Delta Lake};

% ===== Usages =====
\draw[cadre] (16.6,0) rectangle (20.0,-3.2);
\node at (18.3,-1.2) {\logo[2em]{spark}}; \node[et] at (18.3,-2.15) {Requêtes SQL\\(Spark SQL)};
\draw[cadre] (16.6,-3.6) rectangle (20.0,-6.8);
\node at (18.3,-4.8) {\logo[2em]{sklearn}}; \node[et] at (18.3,-5.75) {Prédiction\\des retards};
\draw[cadre] (16.6,-7.2) rectangle (20.0,-10.4);
\node at (18.3,-8.4) {\logo[2em]{streamlit}}; \node[et] at (18.3,-9.35) {Tableau\\de bord};

% ===== Infrastructure =====
\node[fill=lhfonce,text=white,minimum width=20cm,minimum height=0.75cm,anchor=north west] at (0,-10.8)
  {\logo[1.3em]{docker}\textbf{Infrastructure : conteneurs Docker}};

% ===== Flux =====
\draw[flux] (3.2,-3.0) -| (8.1,-4.05) node[pastille,pos=0.3]{1};
\node[et,above] at (5.3,-3.0) {JDBC};
\draw[flux] (3.2,-4.6) -- (4.85,-4.6) node[pastille,midway]{2};
\draw[flux] (5.75,-4.6) -- (7.65,-4.6) node[pastille,midway]{3};
\draw[flux] (8.55,-4.6) -- (10.45,-4.6) node[pastille,midway]{4};
\draw[flux] (10.9,-5.95) -- (10.9,-6.2) node[pastille,right=2pt]{5};
\draw[flux] (11.35,-4.6) -- (13.25,-4.6) node[pastille,midway]{6};
\draw[very thick,lhfonce] (14.15,-4.6) -- (15.55,-4.6);
\draw[flux] (15.55,-4.6) |- (16.6,-1.6) node[pastille,pos=0.75]{7};
\draw[<->,very thick,lhfonce] (15.55,-5.2) -- (16.6,-5.2) node[pastille,midway,below=3pt]{8};
\draw[flux] (15.55,-4.6) |- (16.6,-8.8) node[pastille,pos=0.75]{9};
\end{tikzpicture}}
\caption{Pipeline du TP 1. Tous les services tournent dans des conteneurs Docker.}
\label{c3:fig:pipeline}
\end{figure}

\subsection*{Le pipeline, étape par étape}

\begin{description}[leftmargin=1.4em,labelindent=0pt,itemsep=3pt,font=\normalfont]
\item[\etape{1}] \textbf{Ingestion de la base de vente} (\nb{01\_bronze}) : Spark lit les tables
  \cmd{orders}, \cmd{order\_items} et \cmd{order\_payments} de PostgreSQL par une connexion JDBC et
  les copie en Bronze. Le lot 1 reflète la base au 1\textsuperscript{er} juillet 2018. Le lot 2
  n'apporte que les lignes nouvelles ou modifiées depuis (colonne \cmd{updated\_at}).
\item[\etape{2}] \textbf{Dépôt des fichiers} (\nb{01\_bronze}) : les CSV des référentiels (clients, produits, vendeurs,
  traduction des catégories) sont copiés tels quels dans \cmd{s3://landing/}, un dossier par table et
  par date de dépôt.
\item[\etape{3}] \textbf{Ingestion des fichiers} (\nb{01\_bronze}) : Spark lit ces CSV \emph{en texte},
  sans deviner les types, et les écrit en Bronze. Chaque table Bronze reçoit deux colonnes techniques :
  \cmd{\_source} (d'où vient la ligne) et \cmd{\_ingere\_le} (quand elle a été chargée).
\item[\etape{4}] \textbf{Nettoyage} (\nb{02\_silver}) : chaque table Bronze devient une table
  Silver, avec des colonnes renommées et typées, doublons retirés, catégories traduites, délais de livraison et indicateur de retard
  calculés. Au lot 2, l'instruction \cmd{MERGE} met à jour les commandes existantes et ajoute les
  nouvelles.
\item[\etape{5}] \textbf{Quarantaine} (\nb{02\_silver}) : les lignes qui violent une règle bloquante (commande
  \og livrée \fg{} sans date de livraison, paiement de type inconnu) sont isolées dans des tables de
  quarantaine au lieu d'être supprimées.
\item[\etape{6}] \textbf{Préparation de Gold} (\nb{03\_gold\_sql\_incremental}) : les tables Silver
  sont jointes dans \cmd{gold.commandes}, une ligne par commande, puis agrégées en indicateurs (ventes
  par mois, par catégorie, délais par état).
\item[\etape{7}] \textbf{Analyse} (\nb{03\_gold\_sql\_incremental}) : des requêtes SQL sur Gold répondent aux questions de la direction.
\item[\etape{8}] \textbf{Prédiction} (\nb{04\_ml\_retards}) : un modèle scikit-learn apprend à prévoir les retards à partir
  de \cmd{gold.commandes}, puis écrit ses prédictions dans \cmd{gold.predictions\_retard}.
\item[\etape{9}] \textbf{Restitution} (\nb{dashboard\_tp1.py}) : le tableau de bord Streamlit affiche les indicateurs Gold et
  les commandes au plus fort risque de retard.
\end{description}

La partie 5 (\nb{05\_iceberg\_ou\_delta}) compare ensuite les deux formats de table, Iceberg et Delta Lake.

\begin{center}
\small
\renewcommand{\arraystretch}{1.3}
\begin{tabularx}{\textwidth}{@{}>{\bfseries}p{2.2cm} X X@{}}
\toprule
\normalfont\textbf{Couche} & \textbf{Rôle} & \textbf{Règle d'or} \\
\midrule
Bronze & Copie des sources, avec deux colonnes techniques (\cmd{\_source} et \cmd{\_ingere\_le}) &
  On ne modifie rien. On ajoute, on n'écrase jamais l'historique. \\
Silver & Données nettoyées, typées, dédoublonnées et contrôlées &
  Une ligne invalide n'est pas supprimée en silence, elle va en quarantaine. \\
Gold & Tables prêtes à l'usage (une ligne par commande et des agrégats) &
  Chaque indicateur se recalcule à partir de Silver. \\
\bottomrule
\end{tabularx}
\end{center}

Toutes les tables sont au format \textbf{Apache Iceberg} et enregistrées dans le catalogue sous
le nom \cmd{lakehouse.<couche>.<table>}, par exemple \cmd{lakehouse.silver.orders}.

\subsection*{Un pipeline ELT, de l'OLTP à l'OLAP}

\textbf{ELT et non ETL} : dans un ETL (\emph{Extract, Transform, Load}), les données sont
transformées \emph{avant} d'être chargées, et seule la version propre est conservée. Ce TP suit
l'approche \textbf{ELT} (\emph{Extract, Load, Transform}), celle des Lakehouse. Les données sont
d'abord extraites et chargées telles quelles, puis transformées \emph{dans} le Lakehouse, avec Spark.

\begin{center}
\small
\renewcommand{\arraystretch}{1.25}
\begin{tabularx}{\textwidth}{@{}l l X@{}}
\toprule
\textbf{Lettre} & \textbf{Étapes} & \textbf{Ce qui se passe} \\
\midrule
E -- Extract   & \etape{1} \etape{2} & Lecture de PostgreSQL et dépôt des CSV. \\
L -- Load      & \etape{3}           & Écriture en Bronze, sans aucune transformation. \\
T -- Transform & \etape{4} à \etape{8} & Nettoyage et \cmd{MERGE} en Silver, jointures et agrégats en Gold,
                                       prédictions. \\
\bottomrule
\end{tabularx}
\end{center}

Garder la copie brute en Bronze a deux avantages : si une règle de nettoyage change, on rejoue
Silver et Gold sans relire les sources, et si un chiffre est faux, on remonte jusqu'à la donnée d'origine.

\textbf{De l'OLTP à l'OLAP} : la source PostgreSQL est une base \textbf{OLTP} (\emph{OnLine
Transaction Processing}), celle d'une application qui enregistre des commandes une par une, avec
un modèle \emph{normalisé} (7 tables liées par des clés) et des mises à jour fréquentes. La couche
Gold est \textbf{OLAP} (\emph{OnLine Analytical Processing}). Ses tables sont faites pour lire beaucoup
de lignes et calculer des indicateurs. Bronze et Silver gardent le modèle normalisé de la source,
et c'est en Gold que le modèle devient analytique. Le \emph{format}, lui, est analytique dès Bronze :
Parquet range les données par colonne, et une requête ne lit que les colonnes dont elle a besoin.

% =====================================================================
\section{Déroulé}
% =====================================================================

\textbf{Avant de commencer, copiez les énoncés dans votre dossier personnel.} Dans JupyterLab, ouvrez un
terminal (\emph{File > New > Terminal}) et tapez :
\begin{lstlisting}[style=shell]
cp -r notebooks/tp1 apps/dashboard_tp1.py notebooks/mon_travail/
\end{lstlisting}
Travaillez \textbf{uniquement} dans \nb{notebooks/mon\_travail/}. Git ignore ce dossier, vos fichiers ne
seront donc jamais écrasés, et \cmd{git pull} récupérera sans conflit les mises à jour publiées dans le dépôt.
Ne faites ni \cmd{git commit} ni \cmd{git push} dans le dépôt d'accompagnement.

Les notebooks contiennent les consignes détaillées, et les cellules marquées \cmd{\# À COMPLÉTER} sont à
écrire. Exécutez-les \textbf{dans l'ordre}.
Gardez sous la main l'\textbf{aide-mémoire PySpark} de l'annexe~\ref{ann:aide}, qui contient toutes
les opérations dont vous aurez besoin.


\begin{attention}[Un seul notebook actif à la fois]
chaque notebook ouvert garde sa propre session Spark,
qui réserve jusqu'à 2 Go de mémoire. Avant d'ouvrir le notebook suivant, arrêtez les précédents
(panneau \emph{Running Terminals and Kernels} (barre de gauche) > \emph{Shut Down All}).
Les tables restent dans le Lakehouse.
\end{attention}


% =====================================================================
\section{Partie par partie}
% =====================================================================

\subsection*{Partie 1 -- Bronze \hfill \nb{01\_bronze}}

Vous chargez la base PostgreSQL dans l'état où elle était le \textbf{1\textsuperscript{er} juillet 2018}
(lot 1), vous déposez les fichiers CSV dans la zone \cmd{s3://landing/}, puis vous écrivez sept tables Bronze.

\begin{notion}[On lit les fichiers bruts en texte]\index{couche bronze@couche Bronze}
laisser Spark deviner les types (\cmd{inferSchema})
transforme le code postal \cmd{09790} en nombre \cmd{9790}, et le zéro est perdu pour toujours.
Les types sont fixés plus tard, en Silver, de façon explicite.
\end{notion}

\begin{resultat}
\cmd{SHOW TABLES IN lakehouse.bronze} liste 7 tables. \cmd{bronze.orders} contient environ
86\,600 commandes, car les autres n'existent pas encore au 1\textsuperscript{er} juillet 2018.
\end{resultat}

\capture{01_bronze_notebook}{Notebook \texttt{01\_bronze} : chargement du lot 1 et écriture des tables Bronze.}{c3:fig:nb01}
\capture{01_bronze_resultat}{Résultat attendu de la partie 1 : les 7 tables Bronze et le nombre de commandes.}{c3:fig:res01}

Une table Iceberg n'est, physiquement, qu'un dossier de fichiers dans le stockage S3.
Dans la console (\url{http://localhost:9001}), le dossier \cmd{lakehouse/warehouse/bronze/} contient
un sous-dossier par table, avec ses fichiers de données Parquet (\cmd{data/}) et ses fichiers de
métadonnées (\cmd{metadata/}).

\capture{01_bronze_s3}{Les 7 tables Bronze vues dans la console du stockage S3, avec un dossier par table.}{c3:fig:s3bronze}


\subsection*{Partie 2 -- Silver \hfill \nb{02\_silver}}

Vous mesurez d'abord les défauts (valeurs nulles, doublons), puis vous écrivez une table Silver par
source. Un exemple complet est fourni pour \cmd{customers}, et vous traitez les autres sur le même modèle.

\begin{notion}[Une valeur nulle n'est pas forcément une erreur]\index{valeur nulle}
une commande en cours n'a pas encore de date
de livraison, c'est normal. Une commande \og livrée \fg{} sans date de livraison est une erreur, et
elle part en \textbf{quarantaine}. Entre les deux, une incohérence sans gravité devient une
\textbf{alerte}. On garde la ligne et on la signale.
\end{notion}

\begin{resultat}
le bilan final compare Bronze et Silver table par table. Chaque écart s'explique par une règle
(doublons retirés, lignes en quarantaine).
\end{resultat}

\capture{02_silver_notebook}{Notebook \texttt{02\_silver} : nettoyage, typage et quarantaine.}{c3:fig:nb02}
\capture{02_silver_resultat}{Résultat attendu de la partie 2 : le bilan Bronze / Silver table par table.}{c3:fig:res02}

\subsection*{Partie 3 -- Gold, SQL et ingestion incrémentale \hfill \nb{03\_gold\_sql\_incremental}}

Vous construisez \cmd{gold.commandes} (une ligne par commande) et trois tables d'indicateurs,
puis vous répondez en SQL aux questions de la direction. Vous chargez ensuite le \textbf{lot 2} :
deux mois et demi de nouvelles commandes et de changements de statut.

\begin{notion}[Le modèle de la couche Gold]\index{couche gold@couche Gold}
Gold ne suit pas un schéma en étoile, mais un modèle plus
simple, en deux niveaux (figure~\ref{c3:fig:gold}) :
\begin{itemize}[leftmargin=*,itemsep=1pt,topsep=2pt]
\item \cmd{gold.commandes} est une \textbf{table large dénormalisée} (\emph{One Big Table}), avec une
  ligne par commande, où le client, les montants, les paiements et les délais de livraison sont
  déjà joints. On l'interroge sans aucune jointure.
\item Les tables d'indicateurs (\cmd{ventes\_mensuelles}, \cmd{ventes\_categorie},
  \cmd{livraison\_etat}) sont des \textbf{agrégats} calculés à partir d'elle. Ce sont de petits
  \emph{data marts}, chacun prêt pour un graphique du tableau de bord.
\end{itemize}
Dans un \textbf{schéma en étoile}, on aurait au centre une table de faits (une ligne par article
vendu, avec les montants) reliée par des clés à des tables de dimensions (clients, produits,
vendeurs, dates). Ce modèle évite de répéter les informations et sert de nombreuses analyses, mais
chaque requête demande des jointures. La table large est plus simple à interroger et c'est le
format qu'attend le modèle de ML (une ligne par commande, toutes les variables). Avec un stockage
en colonnes, sa largeur ne coûte presque rien à la lecture. Un modèle dimensionnel, une constellation
de plusieurs tables de faits qui partagent leurs dimensions, sera construit au TP~2.
\end{notion}

\begin{figure}[H]
\centering
\begin{tikzpicture}[
  font=\sffamily\small, >=Stealth,
  tb/.style={draw=#1,thick,rounded corners=3pt,fill=#1!7,align=center,inner sep=4pt,
             font=\sffamily\footnotesize,minimum height=0.8cm},
  lab/.style={font=\sffamily\bfseries\footnotesize,text=black!70}]
\node[lab] at (0,2.4) {Silver (modèle normalisé)};
\node[lab] at (5.2,2.4) {Gold -- table large};
\node[lab] at (10.2,2.4) {Gold -- agrégats};
\foreach \n/\i/\y in {orders/a/1.6, payments/b/0.8, customers/c/0, order\_items/d/-0.8, products/e/-1.6}
  \node[tb=lhsilver,text width=2.6cm] (s\i) at (0,\y) {\texttt{\n}};
\node[tb=lhgold,text width=3.4cm,minimum height=2.4cm] (oc) at (5.2,0)
  {\texttt{\textbf{gold.commandes}}\\[2pt] une ligne par commande\\ client, montants,\\ paiement, délais, retard};
\node[tb=lhgold,text width=3.4cm] (pr) at (10.2,1.6) {\texttt{predictions\_retard}\\ calculée par le modèle de ML};
\node[tb=lhgold,text width=3.4cm] (vm) at (10.2,0.4) {\texttt{ventes\_mensuelles}\\ une ligne par mois};
\node[tb=lhgold,text width=3.4cm] (le) at (10.2,-0.8) {\texttt{livraison\_etat}\\ une ligne par état};
\node[tb=lhgold,text width=3.4cm] (vc) at (10.2,-2.0) {\texttt{ventes\_categorie}\\ une ligne par catégorie};
\foreach \i in {a,b,c,d} \draw[->,thick,black!60] (s\i.east) -- (oc.west);
% ventes_categorie : grain article (order_items + products), et non commande
\draw[thick,black!60] (sd.east) -- (1.75,-0.8) -- (1.75,-2.75);
\draw[thick,black!60] (se.east) -- (1.75,-1.6);
\draw[->,thick,black!60] (1.75,-2.75) -- (8.0,-2.75) |- ([yshift=-5pt]vc.west);
\foreach \t in {pr,vm,le} \draw[->,thick] (oc.east) -- (\t.west);
\draw[->,thick] (oc.east) -- ([yshift=5pt]vc.west);
\end{tikzpicture}
\caption{Le modèle de la couche Gold. Les tables Silver sont jointes une seule fois dans
\texttt{gold.commandes}, et les agrégats et les prédictions sont calculés à partir d'elle
(\texttt{ventes\_categorie} relit aussi les articles et les produits de Silver, car une commande peut couvrir plusieurs catégories).}
\label{c3:fig:gold}
\end{figure}

\begin{notion}[Ingestion incrémentale]\index{ingestion incrementale@ingestion incrémentale}
on ne relit que les lignes modifiées depuis la dernière ingestion,
grâce à la colonne \cmd{updated\_at} (le \emph{filigrane}). L'instruction \cmd{MERGE INTO} applique
ces changements en Silver en une seule transaction (mise à jour des commandes existantes,
insertion des nouvelles). Iceberg garde l'ancienne version de la table, que l'on peut relire
avec \cmd{VERSION AS OF} (\emph{time travel}).
\end{notion}

\begin{resultat}
après le \cmd{MERGE}, \cmd{silver.orders} contient environ 99\,400 commandes, et le nombre de
commandes \cmd{shipped} est divisé par deux environ, car elles ont été livrées entre-temps.
\end{resultat}

\capture{03_gold_notebook}{Notebook \texttt{03\_gold\_sql\_incremental} : tables Gold, requêtes SQL et lot 2.}{c3:fig:nb03}
\capture{03_gold_resultat}{Résultat attendu de la partie 3 : \texttt{silver.orders} après le \texttt{MERGE} et ses deux snapshots.}{c3:fig:res03}


\begin{figure}[H]
\centering
\begin{tikzpicture}[
  font=\sffamily\small, >=Stealth,
  box/.style={draw=#1,thick,rounded corners=3pt,fill=#1!7,text width=4.8cm,align=center,
              inner sep=5pt,font=\sffamily\footnotesize,minimum height=1.25cm},
  row/.style={font=\sffamily\bfseries\footnotesize,anchor=east},
  op/.style={font=\sffamily\scriptsize,text=black!70,align=center}]
\node[font=\sffamily\bfseries] at (0,1.0) {Lot 1 -- 1\textsuperscript{er} juillet 2018};
\node[font=\sffamily\bfseries] at (7.6,1.0) {Lot 2 -- octobre 2018};
\node[row] at (-3.0,0)    {PostgreSQL};
\node[row] at (-3.0,-2.0) {Bronze};
\node[row] at (-3.0,-4.0) {Silver};

\node[box=gray]     (p1) at (0,0)    {toute la table \texttt{ventes.orders}\\ $\approx$ 86\,600 commandes};
\node[box=gray]     (p2) at (7.6,0)  {seulement les lignes où\\ \texttt{updated\_at} > filigrane\\ $\approx$ 14\,500 lignes};
\node[box=lhbronze] (b1) at (0,-2.0) {\texttt{bronze.orders} créée};
\node[box=lhbronze] (b2) at (7.6,-2.0) {lignes \textbf{ajoutées} (\texttt{append})\\ l'historique est conservé};
\node[box=lhsilver] (s1) at (0,-4.0) {\texttt{silver.orders}\\ \textbf{snapshot 1}};
\node[box=lhsilver] (s2) at (7.6,-4.0) {\texttt{MERGE} : $\approx$ 1\,700 mises à jour\\ et $\approx$ 12\,800 insertions\\ \textbf{snapshot 2}};

\draw[->,thick] (p1) -- (b1);
\draw[->,thick] (b1) -- (s1);
\draw[->,thick] (p2) -- (b2);
\draw[->,thick] (b2) -- (s2);
\draw[->,dashed] (b1.east) -- (3.8,-2.0) -- (3.8,0) -- (p2.west);
\node[op,fill=white,inner sep=2pt] at (3.8,-1.0) {filigrane =\\ \texttt{max(updated\_at)}};
\draw[->,dashed] (s2.west) -- node[op,below=2pt]{\emph{time travel} :\\ \texttt{VERSION AS OF}} (s1.east);
\end{tikzpicture}
\caption{Ingestion incrémentale. Au lot 2, on ne relit que les lignes modifiées depuis le lot 1.
Iceberg garde l'ancienne version de la table, que l'on peut toujours interroger.}
\label{c3:fig:incremental}
\end{figure}

\subsection*{Partie 4 -- Prédire les retards \hfill \nb{04\_ml\_retards}}

Vous entraînez deux modèles de classification (scikit-learn) qui estiment, au moment de la commande,
la probabilité d'un retard de livraison, puis vous comparez leurs performances. Les prédictions sont
écrites dans \cmd{gold.predictions\_retard}, que le tableau de bord affichera.

\begin{notion}[Fuite de données]\index{fuite de donnees@fuite de données}
un modèle ne doit utiliser que des informations disponibles au moment de
la prédiction. Le notebook vous fait ajouter volontairement le délai réel de livraison, et
le modèle devient presque parfait\dots{} et totalement inutilisable.
\end{notion}

\begin{resultat}
une AUC ROC modeste, de l'ordre de 0,7, car les causes principales des retards (transporteur, météo,
grèves) ne figurent pas dans les données. Avec la fuite de données, l'AUC dépasse 0,95.
\end{resultat}

\capture{04_ml_notebook}{Notebook \texttt{04\_ml\_retards} : découpage temporel, entraînement et comparaison des modèles.}{c3:fig:nb04}
\capture{04_ml_resultat}{Résultat attendu de la partie 4 : AUC ROC des deux modèles, avec et sans fuite de données.}{c3:fig:res04}

\begin{figure}[H]
\centering
\begin{tikzpicture}[font=\sffamily\small, >=Stealth]
\fill[lhblue!18] (0,0) rectangle (8.4,0.8);
\fill[lhgold!40] (8.4,0) rectangle (13.6,0.8);
\draw[thick] (0,0) rectangle (13.6,0.8);
\draw[thick,dashed] (8.4,-0.35) -- (8.4,1.25);
\node[align=center,font=\sffamily\footnotesize] at (4.2,0.4) {\textbf{apprentissage} : $\approx$ 57\,000 commandes};
\node[align=center,font=\sffamily\footnotesize] at (11.0,0.4) {\textbf{test} : $\approx$ 39\,000};
\node[font=\sffamily\scriptsize,anchor=north] at (0,-0.05)   {sept.\ 2016};
\node[font=\sffamily\scriptsize,anchor=north] at (8.4,-0.35) {1\textsuperscript{er} mars 2018};
\node[font=\sffamily\scriptsize,anchor=north] at (13.6,-0.05) {août 2018};
\draw[->] (0,1.55) -- node[above,font=\sffamily\scriptsize]{temps (date d'achat)} (13.6,1.55);
\end{tikzpicture}
\caption{Découpage temporel : le modèle apprend sur le passé et il est évalué sur des commandes
plus récentes, comme il le serait en production.}
\label{c3:fig:split}
\end{figure}

\subsection*{Partie 5 -- Iceberg ou Delta Lake ? \hfill \nb{05\_iceberg\_ou\_delta}}

Partie courte (30 à 45 minutes), à faire en travail personnel. Toutes les
tables du TP sont au format Iceberg, et Delta Lake, créé par Databricks, est son principal concurrent.
Vous écrivez la même table (\cmd{silver.orders}) dans les deux formats, vous la modifiez, vous
revenez à sa première version, et vous comparez ce que chaque format range dans le stockage.

\begin{notion}[Format de fichier, format de table]\index{format de table}
Parquet est un \emph{format de fichier}. Il range les
données par colonne, mais ne connaît ni versions ni transactions. Iceberg et Delta Lake sont des
\emph{formats de table}. Au-dessus des mêmes fichiers Parquet, ils ajoutent des métadonnées (schéma,
liste des fichiers, historique). Iceberg les range dans un dossier \cmd{metadata/} et passe par un
catalogue, alors que Delta tient un journal de transactions dans \cmd{\_delta\_log/}.
\end{notion}

\begin{resultat}
les deux tables contiennent les mêmes fichiers Parquet. Seules leurs métadonnées diffèrent. Après
l'\cmd{UPDATE}, les anciens fichiers restent dans le stockage (\emph{time travel}), et seule la table
Iceberg apparaît dans le catalogue.
\end{resultat}


% =====================================================================
\section{Le tableau de bord}\label{c3:sec:dashboard}
% =====================================================================

Ouvrez \nb{notebooks/mon\_travail/dashboard\_tp1.py} dans JupyterLab. Le chargement des données et le filtre de période
sont fournis. Complétez les cinq éléments demandés :

\begin{enumerate}[leftmargin=*,nosep]
  \item une ligne de quatre indicateurs : chiffre d'affaires, nombre de commandes, panier moyen, taux de retard ;
  \item la courbe du chiffre d'affaires mensuel ;
  \item le top 10 des catégories ;
  \item le taux de retard par état ;
  \item les commandes au plus fort risque de retard, selon le modèle de la partie 4, avec un seuil réglable.
\end{enumerate}

Pour le lancer, ouvrez un terminal dans JupyterLab (\emph{File > New > Terminal}) :
\begin{lstlisting}[style=shell]
streamlit run notebooks/mon_travail/dashboard_tp1.py --server.address 0.0.0.0
\end{lstlisting}
Puis ouvrez \url{http://localhost:8501} dans votre navigateur. N'utilisez pas l'adresse
\cmd{http://0.0.0.0:8501} affichée dans le terminal, car elle n'est valable qu'à l'intérieur du conteneur.
Chaque enregistrement du fichier met la page à jour.

\capture{dashboard_haut}{Tableau de bord attendu : les quatre indicateurs et le chiffre d'affaires mensuel.}{c3:fig:dash1}
\capture{dashboard_bas}{Tableau de bord attendu : top 10 des catégories, taux de retard par état et commandes à risque.}{c3:fig:dash2}

\begin{attention}
le tableau de bord démarre sa propre session Spark. Sur une machine de 8 Go, arrêtez d'abord le
noyau des notebooks (\emph{Kernel > Shut Down Kernel}) pour libérer la mémoire.
\end{attention}

\begin{astuce}
pensez à l'utilisateur, avec un titre clair, des unités (R\$, \%, jours), et des mois partiels exclus par
défaut. Un graphique qui fait croire à un effondrement des ventes en fin de période est une erreur
d'analyse, pas un détail.
\end{astuce}

% =====================================================================
\section{Synthèse}\label{c3:sec:livrables}
% =====================================================================

Pour faire le point, rédigez une synthèse de ce TP. Elle ne recopie pas le code mais explique
vos choix et vos résultats. Elle couvre :
\begin{enumerate}[leftmargin=*,itemsep=2pt]
  \item \textbf{Architecture} : le schéma de votre pipeline et le rôle de chaque couche.
  \item \textbf{Qualité des données} : les défauts trouvés, les règles appliquées en Silver
        (bloquantes ou simples alertes) et le bilan Bronze / Silver.
  \item \textbf{Ingestion incrémentale} : le principe du filigrane et du \cmd{MERGE}, et ce que montre
        le \emph{time travel}.
  \item \textbf{Réponses à la direction} : les résultats des requêtes SQL, interprétés.
  \item \textbf{Modèle de prédiction} : les variables choisies, les métriques, la comparaison des deux
        modèles et la démonstration de la fuite de données.
  \item \textbf{Iceberg ou Delta Lake} : ce que les deux formats ont en commun, ce qui les distingue,
        et lequel vous choisiriez pour ce Lakehouse, en justifiant.
  \item \textbf{Tableau de bord} : une capture d'écran et deux ou trois enseignements.
  \item \textbf{Réponses aux questions} des notebooks, rédigées et argumentées.
  \item \textbf{Limites et améliorations} : ce que vous feriez avec plus de temps, ou en production.
\end{enumerate}

\begin{aretenir}
  \item La couche Bronze garde les données telles qu'elles sont arrivées, lues en texte, avec leur origine.
  \item La couche Silver type, contrôle et dédoublonne ; les lignes rejetées vont en quarantaine avec leur raison, et une règle peut être bloquante ou une simple alerte.
  \item La couche Gold répond aux questions métier ; Spark SQL interroge les tables comme une base de données.
  \item L'ingestion incrémentale ne charge que les nouvelles lignes, grâce à un filigrane et à \cmd{MERGE} ; chaque écriture crée un instantané que l'on peut relire.
  \item Un modèle ne doit utiliser que des informations connues au moment de la prédiction, sinon ses performances sont trompeuses (fuite de données).
  \item Iceberg et Delta Lake rendent le même service : des transactions et un historique sur des fichiers Parquet.
\end{aretenir}

\begin{plusloin}
Le fonctionnement de \cmd{MERGE} et du voyage dans le temps est décrit dans la spécification d'Iceberg \citep{icebergspec} et dans l'article sur Delta Lake \citep{armbrust2020}. La modélisation des tables Gold sera approfondie au chapitre~\ref{ch:tp2}, et \citet{kimball2013} en reste la référence.
\end{plusloin}
